\documentclass[aspectratio=169]{beamer}
\usepackage{iftex}
\ifPDFTeX
  \usepackage[T1]{fontenc}
  \usepackage{lmodern}
\fi
\usepackage[czech]{babel}
\usepackage{listings}
\usepackage{tikz}
\usetikzlibrary{arrows.meta,positioning,babel}

\newcommand{\CVUTTemplatePath}{../../../utils/presentation-template}
\usepackage{\CVUTTemplatePath/beamerthemeCVUT}
\setbeamertemplate{footline}[frame number]
\beamertemplatenavigationsymbolsempty
\newcommand{\LANG}{CZ}
\newcommand{\FRONTPAGE}{dark}
\title{Knihovny, balíčky a dokumentace}
\subtitle{Základy programování\\Matouš Cejnek\\U12110, FS, ČVUT v Praze}
\author{} \institute{} \date{}
\renewcommand{\headlineA}{} \renewcommand{\headlineB}{} \renewcommand{\headlineC}{}

\definecolor{CodeBlue}{RGB}{38,83,125}
\definecolor{SoftGray}{RGB}{242,244,247}
\definecolor{CodeGreen}{RGB}{38,125,94}
\definecolor{WarningRed}{RGB}{170,45,45}
\setbeamerfont{normal text}{size=\large}
\setbeamerfont{frametitle}{size=\LARGE}
\setbeamercolor{block title}{fg=white,bg=CodeBlue}
\setbeamercolor{block body}{fg=black,bg=SoftGray}
\newcommand{\term}[1]{\textcolor{CodeBlue}{\textbf{#1}}}
\lstdefinestyle{python}{language=Python,basicstyle=\ttfamily\large,
  keywordstyle=\color{CodeBlue}\bfseries,stringstyle=\color{CodeGreen},
  commentstyle=\color{gray},showstringspaces=false,columns=fullflexible,keepspaces=true}
\lstdefinestyle{pythonSmall}{style=python,basicstyle=\ttfamily\normalsize}

\begin{document}

% 0:00--0:01
\begin{frame}
  \clearpage\thispagestyle{empty}\titlepage
\end{frame}

% 0:01--0:04
\begin{frame}{Co si dnes odnesete?}
  \begin{itemize}
    \item Rozlišit modul, balíček, knihovnu a závislost.
    \item Rozdělit vlastní kód do importovatelných souborů a složek.
    \item Navrhnout malé a srozumitelné veřejné rozhraní.
    \item Popsat funkce docstringy a typovými anotacemi.
    \item Najít, nainstalovat a zaznamenat externí knihovnu.
    \item Oddělit knihovní logiku od programu pro uživatele.
  \end{itemize}
\end{frame}

\begin{frame}{Na co navazujeme?}
  \begin{itemize}
    \item Funkce pojmenovávají opakovaný výpočet.
    \item Třída spojuje data a operace nad nimi.
    \item Výjimka oznamuje, že operaci nelze korektně dokončit.
    \item Soubory ukládají data mimo běh programu.
  \end{itemize}
  \begin{block}{Nová otázka}
    Jak uspořádat více funkcí a tříd tak, aby je mohly používat různé programy?
  \end{block}
\end{frame}

\begin{frame}{Kopírování kódu přestává stačit}
  \begin{itemize}
    \item Opravu musíme ručně přenést do každé kopie.
    \item Programy začnou používat různé verze stejné funkce.
    \item Jeden dlouhý soubor se hůře čte, testuje a rozšiřuje.
  \end{itemize}
  \begin{block}{Lepší dohoda}
    Kód uložíme na jednom místě, pojmenujeme jeho veřejné části a ostatní programy je importují.
  \end{block}
\end{frame}

% 0:04--0:09
\begin{frame}{Čtyři základní pojmy}
  \begin{description}
    \item[Modul] jeden importovatelný soubor, například \texttt{core.py}.
    \item[Balíček] složka sdružující související moduly.
    \item[Knihovna] znovupoužitelný kód nabízející určitou funkčnost.
    \item[Závislost] cizí knihovna, kterou náš projekt potřebuje.
  \end{description}
  \vspace{2mm}
  Knihovna může být tvořena jedním modulem i několika balíčky.
\end{frame}

\begin{frame}{Od projektu ke konkrétní funkci}
  \centering
  \begin{tikzpicture}[>=Latex,node distance=4mm,font=\normalsize,
      box/.style={draw,rounded corners,fill=SoftGray,minimum width=28mm,minimum height=11mm}]
    \node[box] (project) {projekt};
    \node[box,right=of project] (package) {balíček};
    \node[box,right=of package] (module) {modul};
    \node[box,right=of module] (object) {funkce / třída};
    \draw[->,thick,CodeBlue] (project)--(package);
    \draw[->,thick,CodeBlue] (package)--(module);
    \draw[->,thick,CodeBlue] (module)--(object);
  \end{tikzpicture}
  \vspace{6mm}

  \texttt{text-project} $\rightarrow$ \texttt{textkit} $\rightarrow$
  \texttt{core} $\rightarrow$ \texttt{count\_words}
\end{frame}

\begin{frame}[fragile]{Nejjednodušší knihovna je modul}
  \texttt{text\_tools.py}
\begin{lstlisting}[style=python]
def count_words(text: str) -> int:
    return len(text.split())
\end{lstlisting}
  \texttt{main.py}
\begin{lstlisting}[style=python]
import text_tools

print(text_tools.count_words("one two three"))
\end{lstlisting}
\end{frame}

\begin{frame}{Co udělá \texttt{import}?}
  \begin{enumerate}
    \item Python modul najde.
    \item Provede jeho příkazy shora dolů.
    \item Vytvoří objekt modulu se jmény funkcí, tříd a proměnných.
    \item Objekt zpřístupní pod jménem použitým v programu.
  \end{enumerate}
  \begin{block}{Praktický důsledek}
    Modul nemá při importu nečekaně spouštět vstup, výpis nebo zápis do souboru.
  \end{block}
\end{frame}

% 0:09--0:14
\begin{frame}[fragile]{Tři běžné podoby importu}
\begin{lstlisting}[style=pythonSmall]
import math
print(math.sqrt(9))

from math import sqrt
print(sqrt(9))

import math as m
print(m.sqrt(9))
\end{lstlisting}
  \texttt{import math} nejzřetelněji ukazuje, odkud jméno pochází. Alias používáme, jen když zlepšuje čitelnost.
\end{frame}

\begin{frame}{Import přidává jména}
  \begin{itemize}
    \item \texttt{math.sqrt} říká: \texttt{sqrt} patří modulu \texttt{math}.
    \item \texttt{from math import sqrt} přidá přímo jméno \texttt{sqrt}.
    \item Dvě importovaná jména se mohou přepsat nebo plést se jménem naším.
  \end{itemize}
  \begin{alertblock}{Raději nepoužívat}
    \texttt{from math import *} skrývá původ jmen a zvyšuje riziko kolizí.
  \end{alertblock}
\end{frame}

\begin{frame}[fragile]{Modul může být i spustitelný program}
\begin{lstlisting}[style=python]
def main() -> None:
    print("Running the application")


if __name__ == "__main__":
    main()
\end{lstlisting}
  Při importu se \texttt{main()} nespustí. Definice lze bezpečně použít z jiného modulu.
\end{frame}

\begin{frame}[fragile]{Když jeden soubor nestačí: balíček}
\begin{lstlisting}[basicstyle=\ttfamily\large,columns=fullflexible]
text-project/
|-- textkit/
|   |-- __init__.py
|   |-- core.py
|   `-- files.py
`-- main.py
\end{lstlisting}
  Moduly dělíme podle odpovědnosti: \texttt{core.py} zpracovává text, \texttt{files.py} pracuje se soubory.
\end{frame}

% 0:14--0:20
\begin{frame}{K čemu je \texttt{\_\_init\_\_.py}?}
  \begin{itemize}
    \item Označuje adresář jako běžný Python balíček.
    \item Spustí se při importu balíčku.
    \item Může zůstat prázdný.
    \item Může vybrat jména, která balíček nabízí uživatelům.
  \end{itemize}
  \begin{block}{Jeho úloha není „vložit sem vše“}
    Má být krátký a především definovat pohodlné veřejné rozhraní.
  \end{block}
\end{frame}

\begin{frame}[fragile]{Veřejné rozhraní balíčku}
  \texttt{textkit/\_\_init\_\_.py}
\begin{lstlisting}[style=python]
from .core import count_words, normalize_whitespace
\end{lstlisting}
  Uživatel potom píše:
\begin{lstlisting}[style=python]
from textkit import count_words

print(count_words("  one   two  "))
\end{lstlisting}
\end{frame}

\begin{frame}{Uživatel má znát rozhraní, ne vnitřnosti}
  \centering
  \begin{tikzpicture}[>=Latex,node distance=10mm,font=\normalsize,
      box/.style={draw,rounded corners,minimum width=38mm,minimum height=13mm}]
    \node[box,fill=SoftGray] (app) {\texttt{main.py}};
    \node[box,fill=CodeBlue!12,right=of app] (api) {\texttt{textkit}};
    \node[box,fill=SoftGray,right=of api] (inside) {\texttt{core.py}, \texttt{files.py}};
    \draw[->,very thick,CodeBlue] (app)--node[above]{veřejné API}(api);
    \draw[->,thick,gray] (api)--node[above]{vnitřek}(inside);
  \end{tikzpicture}
  \vspace{6mm}

  Vnitřní soubory lze přeuspořádat, pokud veřejná jména a jejich chování zůstanou stejné.
\end{frame}

\begin{frame}[fragile]{Tečka znamená import uvnitř balíčku}
  \texttt{textkit/files.py}
\begin{lstlisting}[style=python]
from pathlib import Path

from .core import normalize_whitespace

def load_clean_text(path: Path) -> str:
    text = path.read_text(encoding="utf-8")
    return normalize_whitespace(text)
\end{lstlisting}
  \texttt{.core} znamená modul \texttt{core} ve stejném balíčku.
\end{frame}

\begin{frame}{Kde Python modul hledá?}
  \begin{itemize}
    \item Mezi vestavěnými moduly a ve standardní knihovně.
    \item Ve složkách, které má pro import k dispozici.
    \item Typicky také v projektu, ze kterého program správně spouštíme.
  \end{itemize}
  \begin{block}{Častý problém}
    Vlastní \texttt{random.py}, \texttt{pathlib.py} nebo \texttt{click.py} může zastínit skutečnou knihovnu stejného jména.
  \end{block}
\end{frame}

% 0:20--0:25
\begin{frame}{Knihovna není uživatelské rozhraní}
  \begin{columns}[T]
    \begin{column}{.46\textwidth}
      \begin{block}{Knihovna}
        Přijímá argumenty, vrací hodnoty a při neplatném použití vyvolá konkrétní výjimku.
      \end{block}
    \end{column}
    \begin{column}{.46\textwidth}
      \begin{block}{Aplikace}
        Čte \texttt{input()}, zpracuje argumenty příkazové řádky a vypisuje zprávy uživateli.
      \end{block}
    \end{column}
  \end{columns}
  \vspace{4mm}
  Stejnou knihovnu pak lze použít z terminálu, testu i jiného programu.
\end{frame}

\begin{frame}[fragile]{Neplatné použití oznamujeme výjimkou}
\begin{lstlisting}[style=pythonSmall]
def require_same_size(left: list[float],
                      right: list[float]) -> None:
    if len(left) != len(right):
        raise ValueError("Dimensions must match")
\end{lstlisting}
  \begin{itemize}
    \item Knihovna neví, zda aplikace chce zprávu vypsat, uložit nebo přeložit.
    \item Volající se rozhodne, zda výjimku zachytí.
    \item Chybu hlásíme co nejblíže místu, kde ji rozpoznáme.
  \end{itemize}
\end{frame}

\begin{frame}[fragile]{Speciální metody tvoří přirozené API}
\begin{lstlisting}[style=pythonSmall]
class Vector:
    def __add__(self, other: "Vector") -> "Vector":
        ...

    def __mul__(self, other: "Vector") -> float:
        ...
\end{lstlisting}
  Uživatel může psát \texttt{a + b} a \texttt{a * b}. Operátor má odpovídat obvyklému významu a rozměry kontroluje třída.
\end{frame}

\begin{frame}{Odkud kód pochází?}
  \begin{description}
    \item[Vlastní modul] soubor, který je součástí našeho projektu.
    \item[Standardní knihovna] dodává ji Python, například \texttt{math} nebo \texttt{pathlib}.
    \item[Externí knihovna] instaluje se samostatně, například \texttt{click}.
  \end{description}
  \begin{block}{Závislost projektu}
    Externí knihovna, bez které některá část projektu nefunguje.
  \end{block}
\end{frame}

% 0:25--0:31
\begin{frame}[fragile]{PyPI a \texttt{pip}}
  \begin{itemize}
    \item PyPI je veřejný katalog distribuovaných Python projektů.
    \item \texttt{pip} umí projekt podle jména stáhnout a nainstalovat.
  \end{itemize}
\begin{lstlisting}[basicstyle=\ttfamily\large,columns=fullflexible]
python -m pip install click
\end{lstlisting}
  \begin{block}{Proč \texttt{python -m pip}?}
    Je zřejmé, ke kterému spuštěnému Pythonu použitý \texttt{pip} patří.
  \end{block}
\end{frame}

\begin{frame}[fragile]{Závislosti zapisujeme do souboru}
  \texttt{requirements.txt}
\begin{lstlisting}[basicstyle=\ttfamily\large,columns=fullflexible]
click==8.3.0
\end{lstlisting}
  Instalace záznamu:
\begin{lstlisting}[basicstyle=\ttfamily\large,columns=fullflexible]
python -m pip install -r requirements.txt
\end{lstlisting}
  Soubor dává spolupracovníkům opakovatelný seznam požadavků. Verze omezuje rozdíly mezi instalacemi v různém čase.
\end{frame}

\begin{frame}{Než knihovnu přidáme}
  \begin{enumerate}
    \item Řeší náš problém a má srozumitelnou dokumentaci?
    \item Je udržovaná a funguje s naší verzí Pythonu?
    \item Jaké další závislosti přinese?
    \item Jakou má licenci?
    \item Dokážeme vytvořit nejmenší funkční příklad?
  \end{enumerate}
  \begin{block}{Cena závislosti}
    Ušetřený vlastní kód výměnou za cizí rozhraní, aktualizace a možné změny chování.
  \end{block}
\end{frame}

\begin{frame}{Dokumentace má různé úlohy}
  \begin{description}
    \item[README] proč projekt existuje a jak jej rychle použít.
    \item[Návod] provede konkrétním úkolem krok za krokem.
    \item[Reference] přesně popíše veřejná jména a jejich chování.
    \item[Příklad] ukáže typické použití v krátkém spustitelném kódu.
  \end{description}
  Jediný dlouhý text obvykle nenahradí všechny čtyři potřeby.
\end{frame}

% 0:31--0:38
\begin{frame}[fragile]{Docstring dokumentuje objekt v kódu}
\begin{lstlisting}[style=python]
def count_words(text: str) -> int:
    """Return the number of words in text."""
    return len(text.split())
\end{lstlisting}
  \begin{itemize}
    \item Je to první řetězec uvnitř modulu, třídy nebo funkce.
    \item Popisuje veřejnou dohodu, ne každý krok implementace.
    \item Nástroje jej zobrazí bez otevírání zdrojového souboru.
  \end{itemize}
\end{frame}

\begin{frame}[fragile]{Složitější funkce potřebuje více informací}
\begin{lstlisting}[style=pythonSmall]
def most_common_words(text: str, count: int) -> list[tuple[str, int]]:
    """Return at most count most frequent words.

    Args:
        text: Text to analyze.
        count: Maximum number of returned items.

    Raises:
        ValueError: If count is negative.
    """
    ...
\end{lstlisting}
\end{frame}

\begin{frame}{Co patří do dokumentace funkce?}
  \begin{itemize}
    \item Co funkce dělá a jaký má výsledek.
    \item Význam argumentů a návratové hodnoty.
    \item Výjimky, které má volající očekávat.
    \item Vedlejší účinky: změna objektu, zápis do souboru nebo výpis.
  \end{itemize}
  \begin{block}{Komentář není náhrada názvu}
    Dobré pojmenování a jednoduchý návrh snižují množství nutného vysvětlování.
  \end{block}
\end{frame}

\begin{frame}[fragile]{Dokumentaci zobrazí Python}
\begin{lstlisting}[style=python]
from textkit import count_words

help(count_words)
print(count_words.__doc__)
\end{lstlisting}
  V terminálu lze vytvořit textovou referenci:
\begin{lstlisting}[basicstyle=\ttfamily\large,columns=fullflexible]
python -m pydoc textkit
\end{lstlisting}
\end{frame}

\begin{frame}{Typová anotace je informace, ne zámek}
  \begin{itemize}
    \item Popisuje očekávané typy argumentů a výsledku.
    \item Pomáhá čtenáři, editoru a nástrojům pro statickou kontrolu.
    \item Python ji při běžném spuštění automaticky nevynucuje.
    \item Nenahrazuje kontrolu hodnot, například kladnosti nebo shodných rozměrů.
  \end{itemize}
  \begin{block}{Dvě různé otázky}
    Typ: „Je \texttt{count} celé číslo?“ Hodnota: „Je \texttt{count} nezáporné?“
  \end{block}
\end{frame}

\begin{frame}[fragile]{Anotace zpřesňuje strukturu dat}
\begin{lstlisting}[style=python]
def word_frequencies(text: str) -> dict[str, int]:
    ...

def read_text(path: str) -> str:
    ...

def top_words(text: str) -> list[tuple[str, int]]:
    ...
\end{lstlisting}
  Typ popíše tvar hodnoty; docstring její význam a pravidla použití.
\end{frame}

% 0:38--0:45
\begin{frame}{Časté chyby při návrhu knihovny}
  \begin{itemize}
    \item Modul při importu spustí \texttt{input()} nebo změní soubor.
    \item Uživatel musí importovat z hlubokých vnitřních modulů.
    \item Funkce místo výsledku jen vypisuje na obrazovku.
    \item Veřejná funkce nemá vysvětlené chování při chybě.
    \item Soubor má stejné jméno jako standardní nebo externí knihovna.
    \item Závislost je použita v kódu, ale není zaznamenána.
  \end{itemize}
\end{frame}

\begin{frame}{Rychlá kontrola porozumění}
  \begin{enumerate}
    \item Jaký je rozdíl mezi modulem a balíčkem?
    \item Proč zveřejňovat jména přes \texttt{\_\_init\_\_.py}?
    \item Co znamená tečka v \texttt{from .core import ...}?
    \item Proč knihovna nemá při chybě volat \texttt{print()}?
    \item Co umí typová anotace a co nekontroluje?
    \item Kam zapíšeme externí závislost projektu?
  \end{enumerate}
\end{frame}

\begin{frame}{Co navazuje po přednášce?}
  \begin{columns}[T]
    \begin{column}{.46\textwidth}
      \begin{block}{Seminář: \texttt{vectorlib}}
        Balíček s třídou \texttt{Vector}, sčítáním, skalárním součinem, kontrolou rozměrů a veřejným API.
      \end{block}
    \end{column}
    \begin{column}{.46\textwidth}
      \begin{block}{Cvičení: \texttt{textkit}}
        Rozšiřovaná knihovna pro zpracování textu, soubory, dokumentaci a jednoduché CLI.
      \end{block}
    \end{column}
  \end{columns}
\end{frame}

\begin{frame}{Shrnutí}
  \begin{itemize}
    \item Modul je soubor; balíček sdružuje moduly do jednoho prostoru jmen.
    \item \texttt{\_\_init\_\_.py} může nabídnout malé a stabilní veřejné API.
    \item Knihovna přijímá hodnoty, vrací výsledky a chyby oznamuje výjimkami.
    \item Docstring vysvětluje význam a typová anotace tvar hodnot.
    \item Externí knihovnu vybíráme vědomě a zaznamenáme jako závislost.
    \item Dobrá struktura dovoluje měnit vnitřek bez změny kódu uživatele.
  \end{itemize}
\end{frame}

\end{document}
